%! AP Statistics — Unit 5, Lesson 5.3 — Homework (Answer Key)
\documentclass[10pt]{article}
\usepackage{apstats-article}
\usepackage{apstats-key}
\newcommand{\TallMath}[1]{$\displaystyle #1\rule[-1.4em]{0pt}{3.2em}$}

\begin{document}

\pageheader{Unit 5, Lesson 5.3}{Homework --- Answer Key}
\namedateperiod

\vspace{0.15in}

\begin{enumerate}

  \item Apple weights $N(182, 20)$, $n = 16$.
        \begin{enumerate}[label=\alph*.]
          \item \ans{Random: stated. Independence: $16 \le 10\%$ of all apples --- yes. Normal/Large Sample: population is approximately normal --- yes. All conditions met.}
          \item \ans{$\mu_{\bar{x}} = 182$ g; $\sigma_{\bar{x}} = 20/\sqrt{16} = 5$ g.}
          \item \ans{$z = (190-182)/5 = 1.60$; $P(\bar{x} > 190) = 1 - 0.9452 = 0.0548$. About 5.5\%.}
          \item \ans{$P(\bar{x} > m) = 0.25 \Rightarrow z = 0.674$; $m = 182 + 0.674(5) \approx 185.4$ g. In 25\% of samples of 16 apples, the sample mean weight exceeds 185.4 g.}
        \end{enumerate}

  \item Coffee shop, $n = 36$.
        \begin{enumerate}[label=\alph*.]
          \item \ans{Random: stated. Independence: $36 \le 10\%$ of all hours --- yes. Large sample: $n = 36 \ge 30$ --- yes. Normal model appropriate by CLT.}
          \item \ans{$\sigma_{\bar{x}} = 9/\sqrt{36} = 1.5$. $z_1 = (26-28)/1.5 \approx -1.33$; $z_2 = (31-28)/1.5 = 2.00$; $P = 0.9772 - 0.0918 = 0.8854$. About 88.5\%.}
        \end{enumerate}

  \item Uniform population, $\mu = 100$, $\sigma = 25$.
        \begin{enumerate}[label=\alph*.]
          \item \ansline{No. $n = 4 < 30$ and the population is not normal (it is uniform). The CLT does not guarantee a normal model for $\bar{x}$ with such a small sample.}
          \item \ans{$\sigma_{\bar{x}} = 25/\sqrt{64} = 3.125$. $z = (97-100)/3.125 = -0.96$; $P(\bar{x} < 97) = 0.1685$. About 16.9\%.}
        \end{enumerate}

\end{enumerate}

\begin{extensionbox}
\textbf{Extension.}
\ans{Need $P(|\bar{x}-50| < 2) \ge 0.95$, i.e., $2/\sigma_{\bar{x}} \ge 1.96$. So $\sigma_{\bar{x}} \le 2/1.96 \approx 1.020$. Then $20/\sqrt{n} \le 1.020 \Rightarrow \sqrt{n} \ge 19.61 \Rightarrow n \ge 385$. Need at least $n = 385$.}
\end{extensionbox}

\begin{reflectionbox}
\textbf{Preview:} Lesson 5.4 examines \emph{unbiased estimators} --- statistics whose long-run average equals the population parameter they estimate.
\end{reflectionbox}

\end{document}
